\documentclass[10pt,a4paper]{article}
\usepackage{amsmath,amssymb,amsthm,graphicx}
% Minimal retained preamble for lane-p2 cards taken from
% sources/ahl-16-fixed-points-2019/editorial-preamble.tex
% (doi:10.5802/ahl.16, Annales Henri Lebesgue 2 (2019) 149-185, CC BY 4.0).
%
% Copied verbatim from the editorial preamble: the package list, the
% theorem declarations under the shared counter and the macro block, so that
% every macro and environment of the retained spans is available.
%
% Omitted and why: the class line (the assembler supplies the class); the
% a4paper geometry line and the fontspec block with its explicit lmroman10
% files (encodings and fonts fight the forced xelatex engine, and the
% assembler supplies the page setup).
% Added for this card only: the retained statement is an enumeration whose two
% items are quoted by the article printed numbers 2.2.1 and 2.2.2, both in the
% card's own prose and in the cross-references of the retained proof, so the
% top level of that enumeration is made to display those numbers. The definition
% is global because the enumerate environment redefines its item label locally.
\makeatletter
\gdef\theenumi{2.2.\arabic{enumi}}
\makeatother
\usepackage{amsmath,amssymb,amsthm,mathtools,graphicx,xcolor,hyperref,xurl,xspace,ifthen}
\newtheorem{theo}{Theorem}[section]
\newtheorem{lemm}[theo]{Lemma}
\newtheorem{theorem}[theo]{Theorem}
\newtheorem*{remas*}{Remarks}
\providecommand{\og}{«\nobreakspace}
\providecommand{\fg}{\nobreakspace»}
\newtheorem{lem}[theo]{Lemma}
\newtheorem{lemma}[theo]{Lemma}
\newtheorem{prop}[theo]{Proposition}
\newtheorem{coro}[theo]{Corollary}
\newtheorem{corol}[theo]{Corollary}
\newtheorem{conj}[theo]{Conjecture}
\newtheorem{Proposition}[theo]{Proposition}
\newtheorem{theoreme}[theo]{Theorem}
\theoremstyle{definition}
\newtheorem{defi}[theo]{Definition}
\newtheorem{exam}[theo]{Example}
\newtheorem{exem}[theo]{Example}
\newtheorem{rema}[theo]{Remark}
\newtheorem{nota}[theo]{Notation}
\newtheorem*{rema*}{Remark}
\newtheorem*{defi*}{Definition}
\newtheorem*{theo*}{Theorem}
\newtheorem*{lemm*}{Lemma}
\newtheorem*{prop*}{Proposition}
\newcommand{\enoncename}{}
\newtheorem{corpusenonce}[theo]{\enoncename}
\newenvironment{enonce}[2][]{\renewcommand{\enoncename}{#2}\begin{corpusenonce}}{\end{corpusenonce}}
\newenvironment{enonce*}[2][]{\par\medskip\noindent\textbf{#2.}\itshape}{\par\medskip}
\newenvironment{proofc}[1][\proofname]{\begin{proof}[#1]}{\end{proof}}
\newcommand{\Mc}{\mathcal{M}}
\newcommand{\pointir}{.\ ---\ }
\providecommand{\mathof}[1]{\mathrm{#1}}

\numberwithin{equation}{section}
\numberwithin{figure}{section}
\numberwithin{table}{section}
\def\endappendix{}
%~Mouliné par MaN_auto v.0.21.4 2019-04-30 09:02:09


\usepackage[all]{xy}

\hyphenation{auto-morphism
every
Roelcke
peri-od-ic
peri-od-icy
rep-re-sen-ta-tion
}








\newcommand{\A}{\mathcal A}
\newcommand{\N}{\mathbb{N}}
\newcommand{\Z}{\mathbb{Z}}
\newcommand{\Q}{\mathbb{Q}}
\newcommand{\K}{\mathcal{K}}
\newcommand{\CC}{\mathbb{C}}

\newcommand*{\mA}{\mathbf{A}}
\newcommand*{\mB}{\mathbf{B}}
\newcommand*{\mC}{\mathbf{C}}
\newcommand*{\mD}{\mathbf{D}}
\newcommand*{\mE}{\mathbf{E}}
\newcommand*{\mF}{\mathbf{F}}
\newcommand*{\mH}{\mathbf{H}}
\newcommand*{\mZ}{\mathbf{Z}}
%\newcommand*{\mA}{\mathbf{A}}
%\newcommand*{\mA}{\mathbf{A}}

\newcommand{\restrict}[2]{#1 \upharpoonright #2}

\newcommand{\e}{\varepsilon}

\DeclareMathOperator{\Aut}{Aut}
\DeclareMathOperator{\Stab}{Stab}
\DeclareMathOperator{\Age}{Age}

\newcommand{\actson}{\curvearrowright}

\newcommand{\RUCB}{\mathrm{RUC}_{b}}
\newcommand{\ROB}{\mathrm{Ro}_{b}}
\newcommand{\WAP}{\mathrm{WAP}}
\newcommand{\AP}{\mathrm{AP}}

\graphicspath{{./figures/}}

\newcommand*{\mk}{\mkern -1mu}
\newcommand*{\Mk}{\mkern -2mu}
\newcommand*{\mK}{\mkern 1mu}
\newcommand*{\MK}{\mkern 2mu}

\hypersetup{urlcolor=purple, linkcolor=blue, citecolor=red}

\newcommand*{\romanenumi}{\renewcommand*{\theenumi}{\roman{enumi}}}
\newcommand*{\Romanenumi}{\renewcommand*{\theenumi}{\Roman{enumi}}}
\newcommand*{\alphenumi}{\renewcommand*{\theenumi}{\alph{enumi}}}
\newcommand*{\Alphenumi}{\renewcommand*{\theenumi}{\Alph{enumi}}}
\let\oldexists\exists
\renewcommand*{\exists}{\mathrel{\oldexists}}
\let\oldforall\forall
\renewcommand*{\forall}{\mathrel{\oldforall}}


\title{A fixed point in the orbit closure of a point is equivalent to finite oscillation stability of the functions evaluated at it}
\author{Lionel Nguyen Van Thé}
\date{doi:10.5802/ahl.16; CC BY 4.0}
\begin{document}
\maketitle
\setcounter{section}{2}
\setcounter{enumi}{1}
\setcounter{equation}{10}
\noindent\textit{Editorial note.} The counted unit is the article's Proposition at source lines 233-240 together with its proof at source lines 241-256, a straight adjacent pairing in the source file. The statement environment carries no source label, so the card is keyed by the label of its first enumerated item, 2.2.1, which is the article's own printed number for that item and the handle the rest of the article uses for it. The source is the publisher version of the article in Annales Henri Lebesgue 2 (2019), pages 149-185, DOI 10.5802/ahl.16, released under CC BY 4.0.\par\medskip
\noindent\textit{Editorial note.} The gate was run on the narrowed interval 236-256, which starts at the line of the enumerated item, because the statement's own label sits three lines inside the statement environment and the gate refuses a start more than two lines above it. The card records the complete environment 233-240, which is the only way the enumerated list renders; this is the same statement/label arrangement as the earlier card of this source whose label also sits on an enumerated item.\par\medskip
\noindent\textit{Editorial note.} The licence evidence used here is the observation file kept with the source, license-observation.json, which records the licence as CC BY 4.0, points at the inherited publisher-version evidence and states that the landing page displays no licence line; the underlying evidence files are the same ones the earlier cards of this source rely on. No acceptance record is written anywhere.\par\medskip
\noindent\textit{Editorial note.} The retained proof uses only the numbering of the two enumerated items of its own statement and the definitions of the section: it contains no pointer to another statement and no citation, so this card declares no navigation pointer and needs no bibliography entry. The equation numbers that appear in the proof are the article's own references into this same statement and are displayed as tags by the assembler.\par\medskip
\noindent\textit{Editorial note.} The retained preamble is the source's editorial preamble with three omissions only: the class line, which the assembler supplies; the page geometry; and the fontspec block with its explicit font files, which would fight the forced xelatex engine. Everything else, including the theorem declarations under the shared counter and the whole macro block, is retained verbatim.\par\medskip
\noindent\textit{Editorial note.} The card makes the excerpt display the numbers the article prints for its own objects, because the retained proof quotes them: the top level of the statement's enumeration is redefined in this card's own copy of the preamble to display 2.2.1 and 2.2.2, which are the handles the proof uses, and the section counter is set to two so that the displayed formula of the proof carries the number the article gives it, which the equation counter at ten produces. Revision r002 set those counters but left the enumeration showing its default labels, and revision r003 adds the redefinition so the labels and the cross-references agree. The theorem-like environment of the statement is renumbered by the assembler as 2.1, the first of its section; nothing in the retained span refers to it by number, since the proof quotes only the two enumerated items.\par\medskip
\noindent\textit{Editorial note.} The census package listed this candidate with the statement span 233-239 and the proof span 241-256; the statement environment was re-read from the source and closes at line 240, so the card records 233-240.\par\medskip
\noindent\textit{Editorial note.} rqual sv-mismatch batch fix 2026-10-14 (Chief ruling: re-assemble unpublished only): source\_version aligned verbatim to the source dir license-observation.json sv (package evidence authority); revision bumped for re-assembly. Recipe from lane-s3.\par\medskip
\noindent\textit{Allowed standard prerequisites (editorial).}\begin{itemize}
\item Flows, orbit closures and finite oscillation stability. A flow is an action of a topological group on a compact space, the orbit closure of a point is the closure of its orbit, and a fixed point is a point the group does not move; a finite family of functions is finitely oscillation stable at a point when for every finite set of group elements and every precision there is a point of the orbit on which all the functions of the family vary by less than the precision. The author's own definitions are the ones used by the counted statement and proof.. The counted proposition is the pointwise criterion that identifies the existence of a fixed point in an orbit closure with the finite oscillation stability of the functions evaluated at the point, and the article uses it as the first step of its fixed point theorems.
\end{itemize}\medskip
\begin{prop}
Let $G$ be a topological group, $G\actson X$ a $G$-flow, and $x\in X$. TFAE:
\begin{enumerate}
\item \label{2.2.1} The orbit closure $\overline{G\cdot x}$ contains a fixed point.
\item \label{2.2.2} For every finite $\mathcal F\subset C(X)$, the family $\{ f_{x}:f\in \mathcal F\}$ is finitely oscillation stable.
\end{enumerate}
\end{prop}


\begin{proof}
$\eqref{2.2.1}\Rightarrow \eqref{2.2.2}$: Fix $\mathcal F\subset C(X)$ finite, $H\subset G$ finite, $\varepsilon>0$. Let $y\in \overline{G\cdot x}$ be a fixed point. Thanks to the continuity of the elements of $\mathcal F$, we may find $g\cdot x$ close enough to $y$ so that for every $f\in \mathcal F$ and every $h\in H$, $|f(h\cdot g\cdot x) - f(h\cdot y)|<\e/2$, i.e. $|f(h\cdot g\cdot x) - f(y)|<\e/2$, since $y$ is fixed. Then, for every $f\in \mathcal F$, $h, h'\in H$, we have:
\begin{align*}
\ensuremath
\left|f_{x}(hg) - f_{x}(h'g)\right| & = |f(h\cdot g\cdot x) - f(h'\cdot g\cdot x)| \\
& = |f(h\cdot g\cdot x) - f(y)| + |f(y) - f(h'\cdot g\cdot x)| \\
& < \e.
\end{align*}

$\eqref{2.2.2}\Rightarrow \eqref{2.2.1}$: For $\mathcal F\subset C(X)$ finite, $H\subset G$ finite, $\varepsilon>0$, define
\[
A_{\mathcal F, H, \varepsilon} = \{ y\in \overline{G\cdot x} :\: \forall f \in \mathcal F \ \forall h\in H \ \ |f(h\cdot y) - f(y)|\leq \varepsilon\}.
\]

This defines a family of closed subsets of $\overline{G\cdot x}$. Thanks to~\eqref{2.2.2}, it has the finite intersection property (every finite intersection of its members contains an element of $G\cdot x$). Its intersection is therefore non-empty. Notice now that this intersection consists of fixed points.
\end{proof}

\end{document}
